\documentclass[11pt,a4paper]{article}
\usepackage[utf8]{inputenc}
\usepackage[T1]{fontenc}
\usepackage[polish]{babel}
\usepackage{lmodern}
\usepackage{geometry}
\usepackage[hidelinks]{hyperref}
\geometry{margin=27mm}
\title{Interocepcja, aleksytymia i regulacja emocji}
\author{}
\date{Prospekt badawczy, wersja 0.1.0, 22 listopada 2026}
\begin{document}
\maketitle
\begin{abstract}
Projekt MSC-INTERO-P01 ma status \texttt{IDEA\_REVIEW}. Wyznaczyć związek między samoopisową świadomością sygnałów z ciała, cechami aleksytymii i deklarowaną regulacją emocji z rozdzieleniem interoceptywnej uwagi od trafności interoceptywnej. Dokument zawiera planowane hipotezy i metody, a nie wyniki empiryczne.
\end{abstract}
\section{Problem i cel badania}
Wyznaczyć związek między samoopisową świadomością sygnałów z ciała, cechami aleksytymii i deklarowaną regulacją emocji z rozdzieleniem interoceptywnej uwagi od trafności interoceptywnej.
\section{Hipotezy robocze}
Większe trudności identyfikowania i opisywania uczuć będą związane z bardziej nieadaptacyjnymi wzorcami regulacji emocji, a związek ten może zależeć od wymiaru interocepcji.
\section{Projekt badania}
Przekrojowe badanie ankietowe online osób dorosłych z wykorzystaniem zweryfikowanych wersji językowych MAIA-2, TAS-20 i DERS lub ERQ; przed doborem narzędzi należy zweryfikować licencje, walidacje i strukturę czynnikową w badanej populacji. Warunki uczestnictwa, kryteria wykluczenia, uzasadnienie liczebności próby i procedura świadomej zgody zostaną ustalone przed prerejestracją.
\section{Operacjonalizacja}
Pierwszorzędowy: związek wyniku trudności identyfikowania uczuć z trudnościami regulacji emocji. Drugorzędowe: podskale świadomości interoceptywnej oraz pozostałe wymiary aleksytymii.
\section{Plan analiz}
Modele regresji wielorakiej z oceną współliniowości i rzetelności podskal; eksploracyjna analiza interakcji interocepcja razy aleksytymia. Mediacji przyczynowej nie należy wnioskować na podstawie danych przekrojowych. Przed analizą danych należy zamrozić estymand, reguły czyszczenia danych, plan braków danych i kryteria analiz wrażliwości.
\section{Etyka, bezpieczeństwo i dane}
Nie rozpoczęto rekrutacji. Konieczne jest rozstrzygnięcie wymogów opinii komisji etycznej, minimalizacji danych, retencji i informowania o dobrowolności udziału. Nie zbierać danych pozwalających na identyfikację w publicznym repozytorium.
\section{Wyniki}
Brak danych i wyników. Sekcja zostanie uzupełniona dopiero po zrealizowaniu zatwierdzonego protokołu.
\section{Ograniczenia}
MAIA-2 mierzy samoopisowe aspekty interocepcji, nie obiektywną trafność detekcji sygnałów fizjologicznych. Wspólna metoda pomiaru może zwiększać korelacje.
\section{Dyskusja}
Interpretacja będzie zależała od trafności narzędzi, jakości danych i alternatywnych wyjaśnień. Projekt nie zakłada z góry potwierdzenia hipotez.
\section*{Literatura}
Weryfikacja prac źródłowych i przygotowanie bibliografii z kontrolą DOI należą do kolejnego etapu. Nie umieszczono niezweryfikowanych cytowań.
\end{document}
